\section{Compatible descent maps and the principal signed-data bridge}
\label{sec:sr-compatible-descent}

The geometric construction uses one arithmetic model, its integral
lattice, and its archimedean metric at a time.  We prove two additional
parts of that construction: the finite group descent maps, with their
integral and norm bounds, and the passage from the principal host's
analytic estimates to the signed correction and complementary mean.
The existence of all geometric models with uniform bounds is distinguished
from these implications at the end of the section.

\subsection{The principal host in the coordinates of the signed comparison}
\label{sec:sr-principal-coordinate-bridge}

Put \(D_h=2^{h+2}\), with \(h\geq2\), and
\[
 H=(\mathbb Z/D_h)^\times/\{1,-1\},\qquad n=|H|=2^h .
\]
Let \(q,d\) be distinct odd primes, with \(q\equiv3\pmod8\),
\(d\equiv7\pmod{16}\) and \(q\mid d+1\).  The principal host is
\[
 E=\mathbb Q(\zeta_{D_h})^+(\sqrt{-q},\sqrt{-d}).
\]
The divisibility condition is a property of the selected window pair.
In the frozen arithmetic construction,
\srcid{CharacterBoxes.ExistsPageN.WinPair} includes the field
\srcid{d_q : q \mid d+1}; \srcid{CharacterBoxes.PairGood} retains
a \srcid{WinPair} as its \srcid{pair} field.
The map \srcid{CharacterBoxes.HostOfField.hostOfField} consumes that
same \srcid{PairGood} when making the host.
Thus the principal window selection supplies \(q\mid d+1\) before
the host is passed to the geometric construction.
It is not a consequence of the weaker abstract \srcid{Host.pair}
condition.  Every application of this bridge uses a selected window
host with this retained provenance.

This is the fixed field of the subgroup specified by the dyadic
congruence in the arithmetic construction.  Its real cyclotomic
factor, rather than the full cyclotomic field, is used here.

Write \(\lambda_q,\lambda_d\) for the two quadratic residue
characters.  The cyclotomic Artin map at the modulus \(D_hqd\) is
\[
 a\longmapsto \bigl(\lambda_q(a),\lambda_d(a),[a]_{D_h}\bigr).
\]
The three coordinate maps are simultaneously surjective by the Chinese
remainder theorem.  Their fixed field is the displayed compositum;
in particular \([E:\mathbb Q]=4n\).  If
\(\lambda_q(a)=(-1)^{a_q}\) and
\(\lambda_d(a)=(-1)^{a_d}\), use the additive coordinates
\begin{equation}\label{eq:sr-host-coordinate-swap}
 \operatorname{Gal}(E/\mathbb Q)\ \longrightarrow\
       \Gamma=(\mathbb Z/2)^2\times H,\qquad
 (\lambda_q(a),\lambda_d(a),[a])\longmapsto(a_d,a_q,[a]).
\end{equation}
Thus the first sign in the signed comparison is the \(d\)-sign.
Complex conjugation is \(c=(1,1,0)\), because the dyadic factor is real.
The embeddings carrying the chosen root of \(-d\) to \(i\sqrt d\)
form exactly \(\Phi=\{(0,b,h)\}\).  The characters of
Section~\ref{sec:sr-finite} are consequently
\[
                   \chi=\lambda_d,\qquad \omega=\lambda_q .
\]

For completeness this coordinate choice gives exactly the inertia
conditions of that section.  Outside \(2,q,d\) the extension is
unramified.  At \(d\), cyclotomic inertia maps onto
\(\{0,(1,0,0)\}\).  At \(q\), the field
\(\mathbb Q(\sqrt{-d})\) is split: the congruence
\(-d\equiv1\pmod q\) lifts to a square in \(\mathbb Q_q\).
It is also split at two, because \(-d\equiv1\pmod8\).
For the latter assertion start with an odd root modulo \(8\);
if \(z^2+d=2^j u\), \(j\geq3\), retain \(z\) when \(u\)
is even and replace it by \(z+2^{j-1}\) when \(u\) is odd.
Starting from \(z_3=1\), this defines odd integers \(z_j\) with
\(z_j^2+d\equiv0\pmod{2^j}\) and
\(z_{j+1}-z_j\in\{0,2^{j-1}\}\).
Indeed the change of its square is
\(2^jz_j+2^{2j-2}\equiv2^j\pmod{2^{j+1}}\).
The sequence is Cauchy in \(\mathbb Z_2\), since the valuations of
its successive differences tend to infinity.
Its limit \(z\) satisfies \(z^2+d=0\), by continuity.
Literal compatibility of the intermediate roots modulo \(2^j\)
is not required.
The decomposition groups at \(2,q\) therefore act trivially on the
\(d\)-quadratic subfield.  They lie in \(\Phi\); both their inertia
groups and their Frobenius representatives lie in \(\Phi\).
These are precisely the host conditions used in the signed balance.

Let \(\theta\in\widehat\Gamma\), and let \(\theta^*\) be the primitive
Dirichlet character inducing its cyclotomic Artin character.
Use its actual normalization
\[
                    \ell_\theta=L'(1,\theta^*)/L(1,\theta^*)
                    \quad(\theta\ne1).
\]
Then the inertial coefficient in the signed comparison is exactly
\begin{equation}\label{eq:sr-primitive-inertial-coefficient}
 \frac1{|I_p|}\sum_{u\in I_p}\theta(rF_p+u)
                         =\theta^*(p^r)\qquad(r\geq1).
\end{equation}
Indeed a nontrivial character on \(I_p\) has sum zero, and
the primitive Dirichlet character is ramified at \(p\) in exactly
this case.  If the inertia restriction is trivial, the sum is
\(\theta(F_p)^r\); the primitive character is unramified and its
value at \(p\) is \(\theta(F_p)\).
These assertions follow directly in the cyclotomic realization:
the \(p\)-primary unit factor is the inertia group, a character
kills that factor exactly when its conductor is prime to \(p\),
and Frobenius on the prime-to-\(p\) cyclotomic factor acts by \(p\).
This argument includes primes dividing the larger host modulus but
not the conductor of \(\theta^*\).
Thus all the \(G_\theta\), at both thick and thin nodes, are the
same centered primitive-character ramps as those in the host
analytic estimates.

Here are the estimates needed from that host, stated without a
predicate name.  For the three nontrivial characters of
\(\Gamma/H\), assume
\begin{equation}\label{eq:sr-host-three-quadratics}
 |G_\chi(x)|,\ |G_{\chi\omega}(x)|,\ |G_\omega(x)|
                         \leq x^{1/2+\epsilon}.
\end{equation}
For every \(\theta\in\widehat\Gamma\), assume
\begin{equation}\label{eq:sr-host-thin-analytic}
 |G_\theta(y)|\leq y\exp(-c x^{\epsilon/8})
 \quad\left(y\geq \tfrac13\exp(x^{\epsilon/3})\right),
 \qquad c>0 .
\end{equation}
The first estimate is the principal specialization of the source's
bounded-character estimate, rather than an additional independence
assumption.  To check this, the retained primitive set for the
cyclic subgroup generated by the principal character is \(\{1\}\).
The source first requires triviality on \(H_0\), and here \(H_0=H\).
It then bounds the characters taking value \(-1\) at
\(e=(1,0,0)\), together with those taking value \(1\) there that are
not retained.  The first class contains \(\chi,\chi\omega\);
the second contains \(\omega\), since it is nontrivial.
The surjectivity of the Artin map ensures that none of these
three primitive characters is the principal character.
This unpacks the definitions
\srcid{HostRaw.Bounded}, \srcid{BoundedCharsPointwise} and
\srcid{EAtH.bdG}; it verifies which three bounds are being used.

\begin{srproposition}\label{prop:sr-host-signed-analytic-bridge}
On this principal host, \eqref{eq:sr-host-three-quadratics}
implies the signed correction inequality
\(\mathcal C(x)\leq3x^{1/2+\epsilon}\).
Suppose also that \eqref{eq:sr-host-thin-analytic} holds and that
the thin weights satisfy
\[
 0<w_{u,j}(\rho)\leq2b',\qquad
 b'\exp(x^{\epsilon/3})\leq xb,\qquad
 |a_t/b-x|\leq1,\qquad x\geq3 .
\]
If
\begin{equation}\label{eq:sr-thin-host-budget}
       2f(x+1)\exp(-c x^{\epsilon/8})\leq x^{-1},
\end{equation}
then the complementary mean for these same data satisfies
\(\mathcal R\geq-Q/x\).
\end{srproposition}
\begin{proof}
Taking real parts and applying the triangle inequality to the three
bounds gives the correction inequality.
Put \(X'=\exp(x^{\epsilon/3})\) and
\(\vartheta=\exp(-c x^{\epsilon/8})\).  At a thin node,
\[
 \frac{a_t}{w_{u,j}(\rho)}
   \geq \frac{(x-1)b}{2b'}
   \geq\frac{x-1}{2x}X'\geq X'/3 .
\]
Consequently the modulus of its weighted ramp is at most
\(a_t\vartheta\leq(x+1)b\vartheta\).
In the definition of \(\mathcal R\), for each \(u,t,j\)
there are at most \(f\) summation points, \(2f\) characters and
the prefactor \(1/f\).  Taking absolute values before averaging
over \(u\) gives
\[
 \mathcal R\geq-2f|T|m(x+1)b\vartheta\geq-|T|mb/x .
\]
The last step is \eqref{eq:sr-thin-host-budget}.
Every ramp and every normalization in this calculation belongs to the
one host just fixed, by \eqref{eq:sr-primitive-inertial-coefficient}.
\end{proof}

The threshold in this implication is uniform over the host fields
on the prescribed dyadic size range.  In fact
\(n\leq2x^{2+8\epsilon}\) gives \(f\leq4x^{2+8\epsilon}\).
It is enough to require
\[
       c x^{\epsilon/8}\geq\log16+(4+8\epsilon)\log x .
\]
This holds above a threshold depending on the fixed
\(\epsilon,c\), before choosing any field or node, since
\(x^{\epsilon/8}/\log x\longrightarrow\infty\).

The same observation deals with the repair inequality without a
hidden dependence on the number of nodes.  If
\(f\leq A_f x^{B_f}\), \(\eta\leq A_\eta x^{B_\eta}\),
where the nonnegative exponents and constants are fixed on the
family, and \(b=\exp(x^\epsilon)\), then
\[
 6f x(x+1)(\eta+1)
 \leq12A_f(A_\eta+1)x^{B_f+B_\eta+2}\leq b
\]
eventually, uniformly over that family.
Finally a cutoff can be chosen as the ceiling of the largest of the
finitely many positive ratios \(a_t/w_{u,j}(\rho)\), increased to
one if necessary.  This chooses one cutoff for all the nodes and
changes none of their weights.

\subsection{Integral twisted traces and the exact norm loss}
\label{sec:sr-twisted-trace-descent}

Let \(K\) be a number field, \(R\) a commutative \(K\)-algebra and
\(S\) a finite group of order \(s\).
Suppose \(\rho_g\) are \(K\)-algebra automorphisms satisfying the
right action convention
\[
             \rho_1=\operatorname{id},\qquad
             \rho_{hg}=\rho_g\circ\rho_h .
\]
Fix a character \(\xi:S\to K^\times\), and set
\[
 \operatorname{Tr}_\xi(F)=\sum_{g\in S}\xi(g)^{-1}\rho_g(F),
 \qquad
 A(u)_{h,g}=\rho_{hg}(u),\qquad \Delta(u)=\det A(u).
\]
Assume that multiplication by \(\Delta(u)\) is injective on \(R\).

Let \(V_-,V,V_+\) be finite dimensional \(K\)-subspaces of \(R\).
Assume that \(V\) is \(S\)-stable and that
\[
       V_-\,\rho_h(u)\subset V,\qquad V\,\rho_h(u)\subset V_+
                         \quad(h\in S).
\]
Put \(N=V\cap R_\xi\), where
\(R_\xi=\{F:\rho_g(F)=\xi(g)F\text{ for all }g\}\), and define
\[
             \varphi_h(f)=\operatorname{Tr}_\xi(f\rho_h(u)).
\]

\begin{srproposition}\label{prop:sr-twisted-trace-descent}
The maps \(\varphi_h:V_-\to N\) have zero common kernel, and
\begin{equation}\label{eq:sr-descent-ranks}
           \dim_K V_-\leq s\dim_K N\leq\dim_K V_+ .
\end{equation}
Suppose lattices in the three spaces are preserved by the indicated
actions and products, \(u\) is integral, and the lattice in \(N\)
is \(L_\xi=L\cap R_\xi\), for the chosen lattice \(L\subset V\).
Then every \(\varphi_h\) carries the source lattice into \(L_\xi\).

At each complex embedding let the target norm be Hermitian,
let every \(\rho_g\) act isometrically on \(V\), and assume
\[
                  \|f\rho_h(u)\|_V\leq B\|f\|_{V_-}
                  \quad(h\in S),\qquad B\geq0 .
\]
For these very same maps,
\begin{equation}\label{eq:sr-descent-norms}
                \sum_{h\in S}\|\varphi_h(f)\|_V^2
                              \leq s^3 B^2\|f\|_{V_-}^2 .
\end{equation}
If the integral lattice \(L_\xi\) does not span \(N\), its rank
still satisfies \(s\operatorname{rk}L_\xi\leq\dim_K V_+\).
\end{srproposition}
\begin{proof}
Reindexing \(g\mapsto gk\) gives
\(\rho_k(\operatorname{Tr}_\xi F)=\xi(k)\operatorname{Tr}_\xi F\).
The stability and product assumptions therefore put each
\(\varphi_h(f)\) in \(N\).

If every \(\varphi_h(f)\) is zero, put
\(v_g=\xi(g)^{-1}\rho_g(f)\).  Expanding the trace gives
\[
                    (A(u)v)_h=\varphi_h(f)=0 .
\]
Multiplication by the adjugate gives \(\Delta(u)v=0\).
The assumed injectivity of multiplication by \(\Delta(u)\)
forces \(v=0\); its identity coordinate is \(f\).  Thus
\(f\mapsto(\varphi_h(f))_h\) is an injection \(V_-\to N^S\),
proving the first rank inequality.

For the other inequality use
\(\Psi:N^S\to V_+\), \(\Psi(F)=\sum_h F_h\rho_h(u)\).
If \(\Psi(F)=0\), applying \(\rho_g\) and using the common
eigencharacter of the \(F_h\) yields
\[
                  \sum_h \rho_{hg}(u)F_h=0\qquad(g\in S).
\]
Thus \(A(u)^t F=0\).  Its determinant is again \(\Delta(u)\),
so the same adjugate argument proves that \(\Psi\) is injective.
This proves the second rank inequality.

Every value of \(\xi\) is a root of unity, since \(S\) is finite.
Its inverse is an algebraic integer and a unit, so the
unnormalized trace preserves integral lattices.  The preceding
eigencharacter computation proves integrality in \(L_\xi\),
without dividing by \(s\).

At a complex embedding the same roots of unity have modulus one.
For each \(h\), the triangle inequality, invariance of the target
norm and the multiplication bound give
\[
 \|\varphi_h(f)\|_V
 \leq\sum_g\|\rho_g(f\rho_h(u))\|_V
 \leq sB\|f\|_{V_-}.
\]
Squaring and summing over the \(s\) possible \(h\) proves
\eqref{eq:sr-descent-norms}.  Finally
\(\operatorname{rk}L_\xi=\dim_K(KL_\xi)\leq\dim_K N\),
which proves the last assertion.
\end{proof}

In the geometric application the three spaces are cusp sections at
weights \(G-W,G,G+W\) on one cover, \(S\) is its prescribed level
subgroup, and \(\rho\) is pullback by level changes.
The equality between the deep type lattice and \(L_\xi\), preservation
of the cusp condition, integrality of the pullbacks and isometry of
their Petersson metrics must all refer to this same cover.
Under those identifications the proposition constructs the maps,
rather than merely assuming their existence.
The source names are \srcid{mapsCore},
\srcid{twistedTrace\_equivariant\_and\_injective\_of\_groupDet}
and \srcid{finrank\_le\_card\_mul\_finrank\_eigenspace\_le}.
This is the map input used in the cover-to-covolume argument of
Section~\ref{sec:cov-lattice}.

\subsection{One integral section with all component determinants nonzero}
\label{sec:sr-bounded-determinant-selection}

We next prove the simultaneous integral selection behind the preceding
maps.  Let
\[
                      R=\prod_{Z\in\mathcal Z}R_Z
\]
be a finite product of integral \(K\)-algebras, each embedded in a
field \(\Omega_Z\), and suppose \(S\) permutes the factors through
its action.  Write \(e_Z\) for the component idempotents.
Let \(W\subset R\) be a finite dimensional subspace decomposed as
\(\bigoplus_Z e_ZW\).
For each \(Z\), choose integral elements
\(w_{Z,1},\ldots,w_{Z,n_Z}\) spanning \(e_ZW\) over \(K\).
Suppose there is \(u_0\in W\) for which
\[
                       \Delta(u_0)_Z\ne0\quad(Z\in\mathcal Z).
\]

\begin{srproposition}\label{prop:sr-bounded-determinant-selection}
There are integers \(0\leq m_{Z,i}\leq s^2\) such that
\[
 u=\sum_{Z,i}m_{Z,i}w_{Z,i},\qquad
                       \Delta(u)_Z\ne0\quad(Z\in\mathcal Z).
\]
In particular \(\Delta(u)\) is a nonzero divisor in \(R\).
Suppose \(n_Z\leq N\), and at every archimedean evaluation point
only its one component contributes.
If the nonnegative pointwise norms satisfy the triangle inequality
and every \(w_{Z,i}\) has norm at most \(\beta\Lambda\) on that
component, then \(u\) has pointwise norm at most
\begin{equation}\label{eq:sr-short-determinant-bound}
                              Ns^2\beta\Lambda .
\end{equation}
The bounds are independent of \(|\mathcal Z|\).
\end{srproposition}
\begin{proof}
Consider changing a single coefficient, so the element has the form
\(v+t w_{Z,i}\).  For a component \(Z'\), its group determinant
is a polynomial in \(t\) of degree at most \(s\).
If it is nonzero at the current coefficient, this polynomial
is nonzero, and has at most \(s\) roots in the characteristic-zero
field \(\Omega_{Z'}\).

Only the components in the \(S\)-orbit of \(Z\) can be affected.
Indeed \(\rho_g(w_{Z,i})\) is supported on the component to which
\(\rho_g\) sends \(Z\); all entries in \(A(w_{Z,i})\) vanish at
every other component.  That orbit has at most \(s\) members.
Consequently there are at most \(s^2\) bad integers for this
coefficient.  Among \(0,\ldots,s^2\) choose a good one.
Begin with the coefficients of \(u_0\) in the spanning family and
replace them one at a time.  At every step all component determinants
remain nonzero, so the finite procedure gives the asserted \(u\).
It is integral because its coefficients are ordinary nonnegative
integers.  A product of domains has a nonzero divisor whenever all
of its components are nonzero, proving the next assertion.

At an evaluation point on \(Z\), all generators supported on other
components vanish.  The triangle inequality gives
\[
            \|u\|\leq\sum_{i=1}^{n_Z}m_{Z,i}\|w_{Z,i}\|
                          \leq Ns^2\beta\Lambda ,
\]
which proves the norm estimate.
\end{proof}

The initial nonzero-determinant condition can itself be verified
by orbit evaluations.  Suppose for each component one can find an
evaluation in its function field, together with compatible trivializations
of the section fibers, such that the functionals
\[
                    w\longmapsto \operatorname{ev}_Z(\rho_g w),
                         \qquad g\in S ,
\]
on \(W\) are linearly independent after extending scalars to that
field.  Their joint map is then onto.  Choose values equal to one
at the identity and zero at all other \(g\).
The evaluation of \(A(w)\) is the permutation matrix
\(({\bf1}_{hg=1})_{h,g}\), and has determinant \(1\) or \(-1\).
Thus the determinant polynomial on \(W\) is not identically zero
on that component.

These finitely many nonzero polynomials have a simultaneous
nonvanishing specialization with coefficients in \(K\).
One way to check this even when their coefficients lie in different
extension fields is the following.  For a polynomial with coefficients
in an extension of \(K\), apply a \(K\)-linear functional on the
finite span of its coefficients that is nonzero on one coefficient.
This gives a nonzero polynomial over \(K\); nonvanishing of this
polynomial implies nonvanishing of the original one.
The product of the finitely many resulting polynomials is nonzero.
An infinite field has a point at which a given nonzero polynomial
does not vanish, by induction on its number of variables.
This gives the required \(u_0\).  Thus orbit separation and the
bounded integral generators give one integral section satisfying
all the determinants simultaneously.

Here is the uniform numerical consequence used in the actual descent.
Put \(z\geq1\).  Suppose constants \(C_1,C_3\geq0\), chosen before
all covers, give
\[
 \log s\leq120z^2,\qquad
 N\leq e^{C_3z^3},\qquad \beta\leq e^{C_1z^3},\qquad
                        \Lambda\leq e^{2C_3z^9}.
\]
The same selected section therefore has pointwise norm at most
\[
                         e^{(C_1+3C_3+240)z^9}.
\]
No component count occurs in the constant.
For an invariant Petersson \(L^2\) metric, a pointwise bound \(B\)
for multiplication by \(u\) gives
\(\|f u\|_2\leq B\|f\|_2\), by integration of the squared
pointwise inequality.  Pullback isometry gives this same bound for
every \(\rho_hu\).  Hence
\[
                         C=C_1+3C_3+240
\]
is a uniform exponent in the preceding descent proposition.

At the selected type level the source has the sharper group estimate
\(\log s\leq45z^2\).
If \(b\geq z^{12}\) and \(z\geq2C+200\), the maps satisfy the
required sum-of-squares bound with multiplier \(e^{2b}\):
\[
 \log(s^3e^{2Cz^9})
       \leq135z^2+2Cz^9
       \leq(135+2C)z^9
       \leq z^{10}\leq z^{12}\leq b\leq2b .
\]
This fixes the constants before selecting the geometric cover.
It does not choose a new section separately for each descent map.

For the archive these bounds come from
\srcid{ExistsShortIntegralVectorOfDeterminantBound},
using \srcid{IntegralCuspSlopePowerBound},
\srcid{ShortGeneratorNormComparison},
\srcid{ExistsCuspRegularLevelDeterminant} and
\srcid{CoveringDomainLatticeRegularity}.
The integer coefficient argument is the mathematical content of
\srcid{CardLeOfDetEqZero.exists\_short}.
The proof above supplies the finite selection and descent calculations.
To apply it to every member of the required family one must still
supply the geometric orbit separation, the short spanning families,
and their stated uniform metric comparisons.

\subsection{The simultaneous geometric input still required}
\label{sec:sr-geometric-selection-residual}

The preceding constructions reduce, but do not close, the geometric
existence problem.  Their exact residual inputs are as follows.

First, on every chosen cover the scalar Hodge-weight cusp space must
have the orbit-evaluation separation used above, and its actual
integral lattice must have the short spanning family with
\(N,\beta,\Lambda\) as displayed.
The source constructs orbit separation at its weight
\(c_0\,\mathrm{hodgeWt}\), where the specified \(c_0\) is large enough
for the quotient point-separation theorem.
Its proof uses a descended linearized Hodge line on a finite quotient,
the canonical-bundle identity, nefness and bigness, boundary vanishing
and the applicable point-separation theorem.
Qualitative existence at some larger weight would not prove this
statement at the prescribed weight.

Second, the theta realization, its local factors and the arithmetic
model must be tied to that same cover and those same metrics.
The definite Siegel--Weil theorem is the value of a doubled
Eisenstein series on a specified box section.
In the transposed doubling convention its parameter is
\((d-n)/2\), and its value is the normalized theta average multiplied
by \(1\) when \(n<d\), and by \(2\) when \(d\leq n\).
Meromorphic continuation and this definite value alone do not give
the agreement of two indefinite theta integrals with their ramified
matching maps.
The latter also requires the local intertwining, nonvanishing,
tower and constant-term comparisons, and the extension from
polynomial test vectors to the actual Schwartz data.
The inspected source producer is
\srcid{TheoremKOfMod.theoremK\_head2}; its final assembly
\srcid{SphericalClosure.CloseSph.closeSph} uses precisely these
comparisons.  The new \srcid{ThkBaseNormZeta} theorem gives the
holomorphic value from anisotropic Siegel--Weil and continuation
in the transposed box coordinates; it is one part of this chain.

Third, the arithmetic intersection and cover construction must give
the top inequality \eqref{eq:sr-top} with the four-term error
\eqref{eq:sr-error} for the same selected data.
The descent maps just constructed supply the map and rank part of
the cover-to-covolume argument.  They do not by themselves construct
the deep covers, the degree and rank shares, the component
Hilbert--Samuel comparison, or the height lower bound.
Their source assembly is
\srcid{HodgeBundleWeights.TopCovolG.topCovolG\_alt}, followed by
\srcid{ThirdPlaceSeries.TopCovolSlim2Use.topCovolSlim2\_use\_alt\_alt}.

The selection procedure in the source is more precise than taking
unrelated local choices.  Its
\srcid{nonempty\_hyps\_definite\_side\_of\_closed\_bound}
first takes the maximum of the row-functional threshold \(C_{F,0}\)
and the evaluation-grid threshold \(C_{F,1}\).
It then uses the same extension and the same grid
\(\mathrm{GridA}(\max(C_{F,0},C_{F,1}))\)
for the row functions, local matching, covolume and linked norm bounds.
The later \srcid{nonempty\_hyps\_of\_allRows\_alt}
intersects the finitely many row conditions before selecting its one
family class and its one raw selection.
Its projections provide the spectral data, base rows, height data,
columns and indefinite realization from that raw selection.
This explains how the archive avoids separate incompatible choices.
A standalone construction must prove the component inputs of this
procedure, not just restate its nonempty-package conclusion.

For any finite list of estimates already proved uniformly on such
a selection, taking the maximum of their lower scale thresholds
makes them hold simultaneously.
For constants this argument has to be performed before the family
or host is chosen.  A statement
\(\forall E\,\exists C_E\) cannot supply the fixed \(C\) in
Definition~\ref{def:sr-supply}.
The principal coordinate bridge above does prove that the signed
correction and the complementary mean transfer simultaneously from
one analytically controlled host; the remaining top comparison and
its compatible uniform geometric realization are not established
by this section.

\section{Integral short generators and uniform metric constants}
\label{sec:sr-short-generators}

We prove the arithmetic lattice step in the construction of a short
descent section.  The proof keeps the dependence on the coefficient
field polynomial in its degree and root discriminant.  In particular
it does not require a bound for the index of a chosen short spanning
lattice.

\begin{srlemma}[Quadratic Euclidean transference]
\label{lem:sr-quadratic-transference}
Let \(L\subset\mathbb R^n\) be a full Euclidean lattice, and suppose
every nonzero element of \(L\) has norm at least \(\lambda>0\).
The dual lattice \(L^*\) contains \(n\) linearly independent vectors
of norm at most \(n^2/\lambda\).
\end{srlemma}

\begin{proof}
Minkowski's convex body theorem applied to a cube contained in the
Euclidean ball gives
\(\lambda_1(M)\le\sqrt n\,\det(M)^{1/n}\) for every full lattice \(M\).
Apply this first to \(L\) and then to \(L^*\).  Since
\(\det(L)\det(L^*)=1\), it gives
\[
 \lambda_1(L^*)\le n/\lambda.
\]
Choose a shortest nonzero dual vector and divide out its integral
content to get a primitive \(y_0\in L^*\), still with
\(\|y_0\|\le n/\lambda\).
Primitivity means that there exists \(x_0\in L\) with
\(\langle y_0,x_0\rangle=1\): in a \(\mathbb Z\)-basis of \(L\), the
coordinates of the integral functional \(y_0\) have greatest common
divisor one.

Put \(H=y_0^\perp\) and \(L_H=L\cap H\).
The homomorphism \(L\to\mathbb Z\), \(x\mapsto\langle y_0,x\rangle\),
is onto, so \(L=L_H\oplus\mathbb Zx_0\) as an abelian group.
Thus \(L_H\) is a full lattice in the \((n-1)\)-dimensional space \(H\),
and its shortest nonzero vector has norm at least \(\lambda\).
Inductively choose independent
\(w_1,\ldots,w_{n-1}\in L_H^*\subset H\) with
\(\|w_i\|\le(n-1)^2/\lambda\).
Let
\[
 c_i=\langle w_i,x_0\rangle-
       \operatorname{round}(\langle w_i,x_0\rangle),\qquad
 y_i=w_i-c_i y_0.
\]
Then \(|c_i|\le1/2\), and \(y_i\) is integral on both \(L_H\) and
\(x_0\), so \(y_i\in L^*\).
Its projection to \(H\) is \(w_i\), hence
\(y_0,y_1,\ldots,y_{n-1}\) are independent.  Orthogonality gives
\[
 \|y_i\|^2
 \le\frac{(n-1)^4+n^2/4}{\lambda^2}
 \le\frac{n^4}{\lambda^2}.
\]
Together with the bound for \(y_0\), this completes the induction.
The assertions for \(n=0\) and \(n=1\) follow directly; the displayed
induction is used for \(n\ge2\).
\end{proof}

\begin{srlemma}[Small independent algebraic integers]
\label{lem:sr-small-integers}
Let \(K\) be a number field of degree \(\delta\), discriminant \(D_K\),
and root discriminant \(\operatorname{rd}(K)=|D_K|^{1/\delta}\).
There exist \(\delta\) elements
\(\alpha_1,\ldots,\alpha_\delta\in\mathcal O_K\), independent over
\(\mathbb Q\), such that
\[
 |\sigma(\alpha_i)|\le
 \delta^3\operatorname{rd}(K)
 \quad\hbox{for every embedding }\sigma:K\hookrightarrow\mathbb C.
\]
\end{srlemma}

\begin{proof}
Embed \(K\) in its Minkowski Euclidean space, with norm squared
\(\sum_\sigma|\sigma(x)|^2\), counting both complex embeddings.
The Euclidean dual of the image of \(\mathcal O_K\) is, up to an
orthogonal reflection in the imaginary coordinates, the image of the
codifferent
\(\mathfrak D_K^{-1}=\{y:\operatorname{Tr}_{K/\mathbb Q}
(y\mathcal O_K)\subset\mathbb Z\}\).
This follows by choosing an integral basis and its trace-dual basis:
reflection changes the real Euclidean pairing
\(\operatorname{Re}\sum_\sigma\sigma(x)\overline{\sigma(y)}\)
into \(\operatorname{Tr}(xy)\).

The codifferent has ideal norm \(|D_K|^{-1}\).
For \(0\ne y\in\mathfrak D_K^{-1}\), the inclusion
\(y\mathcal O_K\subset\mathfrak D_K^{-1}\) implies
\(|N_{K/\mathbb Q}y|\ge |D_K|^{-1}\).
The arithmetic--geometric mean inequality therefore gives
\[
 \sum_\sigma|\sigma(y)|^2
 \ge \delta |N_{K/\mathbb Q}y|^{2/\delta}
 \ge \delta\,\operatorname{rd}(K)^{-2}.
\]
In particular every nonzero vector of this dual lattice has norm at
least \(\operatorname{rd}(K)^{-1}\).
Apply Lemma~\ref{lem:sr-quadratic-transference} to the dual lattice.
Its dual, the original integer lattice, has \(\delta\) independent
vectors of Euclidean norm at most
\(\delta^2\operatorname{rd}(K)\).
Each individual embedding is bounded by that norm.
Enlarging \(\delta^2\) to \(\delta^3\) gives the stated form.
\end{proof}

\begin{srproposition}[Short spanning vectors give integral generators]
\label{prop:sr-short-integral-generators}
Let \(V\) be an \(R\)-dimensional vector space over \(K\), and let
\(M\subset V\) be a finitely generated \(\mathcal O_K\)-submodule
spanning \(V\).  For every embedding \(\sigma\), let \(N_\sigma\) be
a seminorm satisfying
\[
 N_\sigma(av)=|\sigma(a)|N_\sigma(v),\qquad
 N_\sigma(v+w)\le N_\sigma(v)+N_\sigma(w).
\]
Suppose \(S\subset M\) spans \(V\) over \(K\) and
\(N_\sigma(s)\le A\) for every \(s\in S\) and every \(\sigma\).
Then \(M\), as an abelian group, is generated by at most \(2\delta R\)
vectors \(w_i\) such that
\begin{equation}\label{eq:sr-generator-bound}
 N_\sigma(w_i)\le
 \max\{1,\delta^4R\operatorname{rd}(K)A\}
 \quad\hbox{for every }i,\sigma.
\end{equation}
\end{srproposition}

\begin{proof}
The case \(R=0\) is immediate.  Choose a \(K\)-basis
\(v_1,\ldots,v_R\) from \(S\), and choose the integers \(\alpha_j\)
of Lemma~\ref{lem:sr-small-integers}.
The \(\delta R\) vectors \(e_{j,k}=\alpha_jv_k\) form a
\(\mathbb Q\)-basis of \(V\); they lie in \(M\), and every seminorm
is at most
\(A'=\delta^3\operatorname{rd}(K)A\).
Write \(D=\delta R\), enumerate this basis by \(e_1,\ldots,e_D\),
and put \(M_0=\sum_\ell\mathbb Ze_\ell\subset M\).

The group \(M\) is free of rank \(D\): it is finitely generated and
torsion free over \(\mathbb Z\), and its rational span is \(V\).
Choose any \(\mathbb Z\)-basis \(b_1,\ldots,b_D\) of \(M\).
No norm bound for this basis is needed.
If \(b_i=\sum_\ell q_{i\ell}e_\ell\), define
\[
 \widetilde b_i=\sum_\ell
       (q_{i\ell}-\lfloor q_{i\ell}\rfloor)e_\ell.
\]
Then \(b_i-\widetilde b_i\in M_0\), so \(\widetilde b_i\in M\).
The \(2D\) vectors \(e_\ell,\widetilde b_i\) generate \(M\), since
they generate every \(b_i\); and
\[
 N_\sigma(\widetilde b_i)
 \le\sum_\ell
       (q_{i\ell}-\lfloor q_{i\ell}\rfloor)N_\sigma(e_\ell)
 \le DA'.
\]
The same bound holds for \(e_\ell\), since \(D\ge1\).
Thus \(DA'=\delta^4R\operatorname{rd}(K)A\) proves
\eqref{eq:sr-generator-bound}.  The index \([M:M_0]\) never enters.
\end{proof}

Here is the precise uniformity consequence used for descent sections.
Let \(z\ge1\) be a size parameter, let \(S_{\rm wt}\ge0\) be the sum
of absolute values of the highest-weight entries, and suppose the
chosen coefficient fields and component cusp lattices satisfy
\begin{align}
 \log\delta&\le391z^2,&
 \log\operatorname{rd}(K)&\le391z^2,\label{eq:sr-field-metric-size}\\
 \log\max(1,R)&\le C_bz^3,&
 \log A&\le C_2(z^3S_{\rm wt}+z^5).\label{eq:sr-span-metric-size}
\end{align}
Here \(C_b,C_2\ge0\) are fixed before fields, components, levels, and
weights are chosen.  Proposition~\ref{prop:sr-short-integral-generators}
gives generators of norm at most \(\Lambda\ge1\), in number \(N\), with
\begin{equation}\label{eq:sr-uniform-short-generators}
 \log\Lambda\le C_3(z^3S_{\rm wt}+z^5),\qquad
 \log\max(1,N)\le C_3z^3,\qquad
 C_3=2000+C_b+C_2.
\end{equation}
Indeed the first logarithm is bounded by the maximum of zero and
\[
 4\log\delta+\log R+\log\operatorname{rd}(K)+\log A
 \le1955z^2+C_bz^3+C_2(z^3S_{\rm wt}+z^5),
\]
which is at most the first right side of
\eqref{eq:sr-uniform-short-generators}; and
\(\log N\le\log2+\log\delta+\log R\le(392+C_b)z^3\).
For a zero component space take the empty generating family and
\(\Lambda=1\).

The component sup norm and rank bounds have a common Hilbert-space
normalization.  Let \(\mathcal H\) be a finite-dimensional Hilbert
space of line-bundle sections on a component with a probability measure \(\mu\),
and assume its norm is the corresponding \(L^2\) norm.
If point evaluation has squared norm at most \(\beta\), then
\[
 |f(y)|^2\le\beta\|f\|_2^2,\qquad
 \dim\mathcal H\le\beta.
\]
To prove the second assertion, choose an orthonormal basis \(f_i\).
The evaluation estimate is equivalent to
\(\sum_i|f_i(y)|^2\le\beta\).
Integrating and using \(\mu(X)=1\) gives
\(\dim\mathcal H=\sum_i\int|f_i|^2\le\beta\).
Thus, if
\begin{equation}\label{eq:sr-bergman-input}
 1\le\beta\le e^{C_1z^3},
\end{equation}
both the pointwise factor \(\sqrt\beta\le\beta\) and the logarithm of
the component rank are controlled by this one fixed \(C_1\).
The measure must be the same normalized component measure used in
the integral lattice metric; an unnormalized quotient measure would
insert its volume into the rank inequality.

For completeness we now combine these bounds with the simultaneous
group determinant construction.  Suppose the prescribed scalar
Hodge weight satisfies \(S_{\rm wt}\le z^5\), the finite level group
has cardinal \(s\) with \(\log s\le120z^2\), and a cusp section with
nonzero group determinant exists on every component.
Apply the componentwise integer-coefficient selection proved in
Section~\ref{sec:sr-compatible-descent}.  It uses coefficients in
\(\{0,\ldots,s^2\}\) in a single simultaneous choice.
The resulting integral section has pointwise norm at most
\(N s^2\beta\Lambda\), hence multiplication by any of its level
translates has norm at most \(B\), where
\begin{equation}\label{eq:sr-uniform-multiplier}
 \log B\le (C_1+3C_3+240)z^9.
\end{equation}
This follows from
\(\log N\le C_3z^3\),
\(\log\Lambda\le2C_3z^9\),
\(\log\beta\le C_1z^3\), and \(2\log s\le240z^2\).
Every constant in \eqref{eq:sr-uniform-multiplier} is chosen before
the level group or component is selected.

The integral maps
\(\varphi_h(f)=\sum_g\xi(g)^{-1}\rho_g(f\rho_hu)\)
therefore satisfy
\[
 \sum_h\|\varphi_h(f)\|^2\le s^3 B^2\|f\|^2.
\]
This uses only the isometry of level translations and the multiplier
bound just proved.  There is no division by \(s\) in the maps, so
their integrality is retained at primes dividing \(s\).

These arguments close the passage from a uniformly short integral
\(K\)-spanning family to the simultaneous integral descent maps.
For an actual Shimura family the remaining input
\eqref{eq:sr-span-metric-size} must concern the prescribed algebraic
theta sections, their common integral lattice, and the same
Petersson metric.  Abstract finite generation does not prove it.

More precisely, the frozen construction supplies this input in the
following order.  Its theta spanning assertion uses the identical
restricted tensor model, the integral Fourier datum
\(\widehat{\mathbf1}_{x_0}\), unimodular positive definite auxiliary
lattices trivial at the prescribed bad places, and integral harmonic
polynomials.  One denominator \(N_a\), valid for every section in that
theta family at the chosen weight, has
\(\log N_a\le C_{\rm int}z^2\kappa\); its raw pointwise estimates have
exponent \(C_{\rm raw}z^2\kappa\).  Multiplication by a reflex-ideal
scalar balances the lattice index at every embedding without changing
the \(K\)-span.  These are the substantive inputs to
\srcid{ExistsGoodCoverAtMinimalLevel}; the theorem name alone is not
a proof of them.
To pass from these even theta weights to a general dominant weight of
gap at least \(8fr+10\), one takes auxiliary even rank
\(8(fr+1)\), a key-unit twisting section, and a flat-line twisting
section with nonnegative finite valuations.  Their weights must add
exactly to the desired weight.  Their archimedean factors have
logarithms bounded by the key-unit term and
\[
 (|c_1|+q\log\operatorname{rd}(E))\sum_\tau|n_{0,\tau}|
 +\sum_{p\ {\rm bad}}\log p
 +\tfrac12\log\operatorname{rd}(K)+1,
 \qquad |c_1|\le9rL.
\]
Together with the proved bad-prime bound and the weight-size
inequalities, the source gives
\(C_2=11C_{\rm even}+C_{\rm bad}+500\) in
\eqref{eq:sr-span-metric-size}.
Those theta spanning, integral-denominator, key-unit, flat-line, and
Bergman constructions must all be established with these
normalizations to discharge the geometric input; the lattice and
constant arguments above do not replace them.
